% GENERATED FILE. Edit canonical lessons, then run npm run generate:books.
% canonical-source-hash: fnv1a64:f12f1dd4238aba84
% canonical-lessons: TE-C328-toluta, TE-C328-anantaram, TE-C328-ettakelaku, TE-C328-marusati-roju, TE-C328-gurtucesuko

\chapter{First, Then, At Last}
\label{ch:te-first-then-at-last}
\hlchaptermodality{328}

\begin{chapteropening}
\textbf{By the end of this chapter:} \emph{I can say at first, afterwards, at long last, the next day and to recall.}

Complete the last lesson of chapter 328: I can say at first, afterwards, at long last, the next day and to recall.
\end{chapteropening}

\section[\texorpdfstring{\te{తొలుత} (toluta)}{తొలుత (toluta)}]{\te{తొలుత} (toluta) — at first}

\label{lesson:TE-C328-toluta}



\begin{quote}
Before the new one: say the Telugu for a kingdom, then the Telugu for to reign over.
\end{quote}

\subsection*{You'll want to know: \te{తొలుత}}
\textbf{\te{తొలుత}} — \emph{toluta} — \textquotedblleft{}at first\textquotedblright{}.

\begin{quote}
\te{తొలుత} \te{నాకు} \te{భయం} \te{వేసింది}.
\end{quote}

\begin{quote}
\emph{toluta nāku bhayaṁ vēsindi.}
\end{quote}

\begin{quote}
At first I felt afraid.
\end{quote}

A storytelling word, a little formal; in speech \te{మొదట} says the same.

\begin{cousinweb}
\textbf{\te{తొలుత}} (\emph{toluta}) is native Dravidian, built on \emph{toli}, \textquotedblleft{}beginning; first, former\textquotedblright{}, which DEDR 3516 lists with Tamil \emph{tol}, \textquotedblleft{}old, ancient\textquotedblright{}.
\end{cousinweb}

\begin{grammarlens}[title={What you've built}]
One more word for telling what happened.
\end{grammarlens}

\subsection*{Your turn}
\begin{itemize}
\raggedright
  \item \emph{Say it:} \emph{toluta}
  \item \emph{Say it:} \emph{toluta}, once more
  \item \emph{Say it:} \emph{toluta nāku bhayaṁ vēsindi.} — At first I felt afraid.
  \item \emph{Recall:} say the Telugu for a kingdom, then the Telugu for to reign over, then say \emph{toluta} again
\end{itemize}

\begin{tcolorbox}[breakable,colback=teal!4,colframe=teal!35!black,title={Before you move on}]
What does \te{తొలుత} mean? (\textquotedblleft{}At first\textquotedblright{}.) Say it once more.
\end{tcolorbox}
\section[\texorpdfstring{\te{అనంతరం} (anantaraṁ)}{అనంతరం (anantaraṁ)}]{\te{అనంతరం} (anantaraṁ) — afterwards; after that}

\label{lesson:TE-C328-anantaram}



\begin{quote}
Before the new one: say the Telugu for to reign over, then the Telugu for at first.
\end{quote}

\subsection*{You'll want to know: \te{అనంతరం}}
\textbf{\te{అనంతరం}} — \emph{anantaraṁ} — \textquotedblleft{}afterwards; after that\textquotedblright{}.

\begin{quote}
\te{భోజనం} \te{అనంతరం} \te{అందరూ} \te{ఇంటికి} \te{వెళ్ళారు}.
\end{quote}

\begin{quote}
\emph{bhōjanaṁ anantaraṁ andarū iṇṭiki veḷḷāru.}
\end{quote}

\begin{quote}
After the meal everyone went home.
\end{quote}

It follows a noun as \te{తరువాత} does, but sounds more formal: news readers use it.

\begin{cousinweb}
\textbf{\te{అనంతరం}} (\emph{anantaraṁ}) is borrowed from Sanskrit \emph{anantaram}, \textquotedblleft{}immediately after\textquotedblright{}, from \emph{an-antara}, \textquotedblleft{}having no interval\textquotedblright{} (CDIAL 287).
\end{cousinweb}

\begin{grammarlens}[title={What you've built}]
One more word for telling what happened.
\end{grammarlens}

\subsection*{Your turn}
\begin{itemize}
\raggedright
  \item \emph{Say it:} \emph{anantaraṁ}
  \item \emph{Say it:} \emph{anantaraṁ}, once more
  \item \emph{Say it:} \emph{bhōjanaṁ anantaraṁ andarū iṇṭiki veḷḷāru.} — After the meal everyone went home.
  \item \emph{Recall:} say the Telugu for to reign over, then the Telugu for at first, then say \emph{anantaraṁ} again
\end{itemize}

\begin{tcolorbox}[breakable,colback=teal!4,colframe=teal!35!black,title={Before you move on}]
What does \te{అనంతరం} mean? (\textquotedblleft{}Afterwards; after that\textquotedblright{}.) Say it once more.
\end{tcolorbox}
\section[\texorpdfstring{\te{ఎట్టకేలకు} (eṭṭakēlaku)}{ఎట్టకేలకు (eṭṭakēlaku)}]{\te{ఎట్టకేలకు} (eṭṭakēlaku) — at long last}

\label{lesson:TE-C328-ettakelaku}



\begin{quote}
Before the new one: say the Telugu for at first, then the Telugu for afterwards.
\end{quote}

\subsection*{You'll want to know: \te{ఎట్టకేలకు}}
\textbf{\te{ఎట్టకేలకు}} — \emph{eṭṭakēlaku} — \textquotedblleft{}at long last\textquotedblright{}.

\begin{quote}
\te{ఎట్టకేలకు} \te{బస్సు} \te{వచ్చింది}.
\end{quote}

\begin{quote}
\emph{eṭṭakēlaku bassu vaccindi.}
\end{quote}

\begin{quote}
At long last the bus came.
\end{quote}

It tells you the wait was long; \te{చివరికి} only says what came last.

\begin{cousinweb}
\textbf{\te{ఎట్టకేలకు}} (\emph{eṭṭakēlaku}) is a Telugu adverb of time; its origin is not settled, and DEDR has no entry for it.
\end{cousinweb}

\begin{grammarlens}[title={What you've built}]
One more word for telling what happened.
\end{grammarlens}

\subsection*{Your turn}
\begin{itemize}
\raggedright
  \item \emph{Say it:} \emph{eṭṭakēlaku}
  \item \emph{Say it:} \emph{eṭṭakēlaku}, once more
  \item \emph{Say it:} \emph{eṭṭakēlaku bassu vaccindi.} — At long last the bus came.
  \item \emph{Recall:} say the Telugu for at first, then the Telugu for afterwards, then say \emph{eṭṭakēlaku} again
\end{itemize}

\begin{tcolorbox}[breakable,colback=teal!4,colframe=teal!35!black,title={Before you move on}]
What does \te{ఎట్టకేలకు} mean? (\textquotedblleft{}At long last\textquotedblright{}.) Say it once more.
\end{tcolorbox}
\section[\texorpdfstring{\te{మరుసటి} \te{రోజు} (marusaṭi rōju)}{మరుసటి రోజు (marusaṭi rōju)}]{\te{మరుసటి} \te{రోజు} (marusaṭi rōju) — the next day}

\label{lesson:TE-C328-marusati-roju}



\begin{quote}
Before the new one: say the Telugu for afterwards, then the Telugu for at long last.
\end{quote}

\subsection*{You'll want to know: \te{మరుసటి} \te{రోజు}}
\textbf{\te{మరుసటి} \te{రోజు}} — \emph{marusaṭi rōju} — \textquotedblleft{}the next day\textquotedblright{}.

\begin{quote}
\te{మరుసటి} \te{రోజు} \te{అతను} \te{మళ్ళీ} \te{వచ్చాడు}.
\end{quote}

\begin{quote}
\emph{marusaṭi rōju atanu maḷḷī vaccāḍu.}
\end{quote}

\begin{quote}
The next day he came again.
\end{quote}

\te{మరుసటి} works with other time words too: \te{మరుసటి} \te{వారం}, the next week.

\begin{cousinweb}
\textbf{\te{మరుసటి} \te{రోజు}} (\emph{marusaṭi rōju}) is a Telugu phrase: \emph{marusaṭi}, \textquotedblleft{}the following\textquotedblright{}, is native Dravidian (DEDR 4766, beside Tamil \emph{maṟu}, \textquotedblleft{}other, next\textquotedblright{}), and \emph{rōju} is the everyday word for a day.
\end{cousinweb}

\begin{grammarlens}[title={What you've built}]
One more word for telling what happened.
\end{grammarlens}

\subsection*{Your turn}
\begin{itemize}
\raggedright
  \item \emph{Say it:} \emph{marusaṭi rōju}
  \item \emph{Say it:} \emph{marusaṭi rōju}, once more
  \item \emph{Say it:} \emph{marusaṭi rōju atanu maḷḷī vaccāḍu.} — The next day he came again.
  \item \emph{Recall:} say the Telugu for afterwards, then the Telugu for at long last, then say \emph{marusaṭi rōju} again
\end{itemize}

\begin{tcolorbox}[breakable,colback=teal!4,colframe=teal!35!black,title={Before you move on}]
What does \te{మరుసటి} \te{రోజు} mean? (\textquotedblleft{}The next day\textquotedblright{}.) Say it once more.
\end{tcolorbox}
\section[\texorpdfstring{\te{గుర్తుచేసుకో} (gurtucēsukō)}{గుర్తుచేసుకో (gurtucēsukō)}]{\te{గుర్తుచేసుకో} (gurtucēsukō) — to recall; to remember}

\label{lesson:TE-C328-gurtucesuko}



\begin{quote}
Before the new one: say the Telugu for at long last, then the Telugu for the next day.
\end{quote}

\subsection*{You'll want to know: \te{గుర్తుచేసుకో}}
\textbf{\te{గుర్తుచేసుకో}} — \emph{gurtucēsukō} — \textquotedblleft{}to recall; to remember\textquotedblright{}.

\begin{quote}
\te{ఆమె} \te{తన} \te{బాల్యం} \te{గుర్తుచేసుకుంది}.
\end{quote}

\begin{quote}
\emph{āme tana bālyaṁ gurtucēsukundi.}
\end{quote}

\begin{quote}
She recalled her childhood.
\end{quote}

Past (I …): \te{గుర్తుచేసుకున్నాను}, \emph{gurtucēsukunnānu}.

Built like \te{గుర్తుపెట్టుకో}: \te{గుర్తు}, a mark in the mind, with \te{చేసుకో}, to make for oneself.

\begin{cousinweb}
\textbf{\te{గుర్తుచేసుకో}} (\emph{gurtucēsukō}) is a native Telugu compound: \emph{gurtu}, \textquotedblleft{}a sign, a mark in the memory\textquotedblright{} (DEDR 1847), and \emph{cēsukō}, \textquotedblleft{}to do for oneself\textquotedblright{}, the -\te{కో} form of \emph{cēyu}.
\end{cousinweb}

\begin{grammarlens}[title={What you've built}]
One more word for telling what happened.
\end{grammarlens}

\subsection*{Your turn}
\begin{itemize}
\raggedright
  \item \emph{Say it:} \emph{gurtucēsukō}
  \item \emph{Say it:} \emph{gurtucēsukō}, once more
  \item \emph{Say it:} \emph{āme tana bālyaṁ gurtucēsukundi.} — She recalled her childhood.
  \item \emph{Recall:} say the Telugu for at long last, then the Telugu for the next day, then say \emph{gurtucēsukō} again
\end{itemize}

\begin{tcolorbox}[breakable,colback=teal!4,colframe=teal!35!black,title={Before you move on}]
What does \te{గుర్తుచేసుకో} mean? (\textquotedblleft{}To recall; to remember\textquotedblright{}.) Say it once more.
\end{tcolorbox}
